\documentclass[12pt,oneside]{report}

% ------------------------------------------------
% PAQUETES
% ------------------------------------------------
\usepackage[spanish]{babel}
\usepackage[utf8]{inputenc}
\usepackage[T1]{fontenc}
\usepackage{geometry}
\usepackage{setspace}
\usepackage{graphicx}
\usepackage{titlesec}
\usepackage{tocbibind}
\usepackage{hyperref}
\usepackage{cite}

\geometry{left=3cm,right=2.5cm,top=3cm,bottom=2.5cm}
\onehalfspacing

% ------------------------------------------------
% DATOS DE LA TESIS (MODIFICA AQUÍ)
% ------------------------------------------------
\newcommand{\titulo}{Cyber-Escudo: Plataforma Web de Denuncia con Clasificación Inteligente y Seguimiento Seguro}
\newcommand{\autor}{José Jiménez}
\newcommand{\carrera}{Ingeniería en Sistemas Computacionales}
\newcommand{\director}{Nombre del Director de Tesis}
\newcommand{\unidad}{Unidad Académica (Ej. ESCOM)}
\newcommand{\anio}{2026}

% ------------------------------------------------
\begin{document}

% ------------------------------------------------
% PORTADA
% ------------------------------------------------
\begin{titlepage}
\begin{center}

{\Large INSTITUTO POLITÉCNICO NACIONAL} \\[0.5cm]
{\large \unidad} \\[1.5cm]

{\Large \textbf{\titulo}} \\[1.5cm]

TESIS \\[0.3cm]
QUE PARA OBTENER EL TÍTULO DE \\[0.3cm]
{\textbf{\carrera}} \\[1cm]

PRESENTA: \\[0.3cm]
{\textbf{\autor}} \\[1cm]

DIRECTOR DE TESIS: \\[0.3cm]
{\director} \\[2cm]

Ciudad de México \\[0.3cm]
\anio

\end{center}
\end{titlepage}

% ------------------------------------------------
% PRELIMINARES (ROMANOS)
% ------------------------------------------------
\pagenumbering{roman}

\chapter*{Resumen}
\addcontentsline{toc}{chapter}{Resumen}

En la actualidad, la integración de la tecnología en la vida cotidiana ha generado un entorno digital en el que coexisten oportunidades y riesgos. Fenómenos como el ciberacoso, la suplantación de identidad y la violencia digital representan una problemática social emergente que afecta a estudiantes y docentes.

El presente trabajo propone el desarrollo de la plataforma \textbf{Cyber-Escudo}, orientada a proporcionar un canal digital seguro para el registro anónimo de denuncias, su clasificación mediante inteligencia artificial y el seguimiento transparente de los casos, garantizando protección de datos personales y cumplimiento de la normativa institucional del Instituto Politécnico Nacional.

\tableofcontents
\newpage

% ------------------------------------------------
% CONTENIDO PRINCIPAL
% ------------------------------------------------
\pagenumbering{arabic}

% ------------------------------------------------
\chapter{Introducción}

En la última década, México ha experimentado una acelerada adopción tecnológica. De acuerdo con el Instituto Nacional de Estadística y Geografía (INEGI), más de 100 millones de personas utilizan internet en el país \cite{ref1}. Sin embargo, este crecimiento ha incrementado los riesgos asociados al entorno digital.

Según el Módulo sobre Ciberacoso (MOCIBA), millones de personas han sido víctimas de agresiones digitales \cite{ref2}. Esta problemática también impacta al ámbito educativo, afectando el desempeño académico y el bienestar emocional.

Existe una brecha significativa entre víctimas y denuncias formales, derivada de la desconfianza institucional, temor a represalias y desconocimiento de los mecanismos de reporte.

% ------------------------------------------------
\chapter{Planteamiento del Problema}

En el Instituto Politécnico Nacional, la Defensoría de los Derechos Politécnicos es responsable de proteger los derechos de la comunidad. No obstante, enfrenta limitaciones operativas tales como:

\begin{itemize}
    \item Saturación administrativa.
    \item Procesos manuales y poco digitalizados.
    \item Demoras en el seguimiento de casos.
    \item Falta de herramientas tecnológicas integradas.
\end{itemize}

Estas deficiencias impactan la percepción de confianza y eficiencia institucional.

% ------------------------------------------------
\chapter{Objetivos}

\section{Objetivo General}

Desarrollar una plataforma web de denuncia que integre registro seguro, clasificación automatizada mediante inteligencia artificial y seguimiento transparente de los casos.

\section{Objetivos Específicos}

\begin{itemize}
    \item Analizar el estado del arte de plataformas de denuncia.
    \item Diseñar la arquitectura del sistema con enfoque en seguridad y escalabilidad.
    \item Diseñar un modelo de base de datos estructurado.
    \item Implementar mecanismos de anonimato y cifrado.
    \item Desarrollar módulo de seguimiento por folio.
\end{itemize}

% ------------------------------------------------
\chapter{Justificación}

La falta de canales accesibles y privados limita la denuncia de casos de violencia digital. La implementación de una plataforma tecnológica permitirá optimizar tiempos de respuesta, fortalecer la transparencia y proteger la identidad del denunciante.

La complejidad técnica radica en la implementación de algoritmos de clasificación supervisada y técnicas de procesamiento de lenguaje natural, así como en la aplicación de medidas de seguridad como cifrado y control de accesos.

% ------------------------------------------------
\chapter{Estado del Arte}

Se analizarán plataformas de denuncia implementadas por dependencias gubernamentales y universidades públicas, evaluando criterios como:

\begin{itemize}
    \item Nivel de anonimato.
    \item Seguridad de la información.
    \item Automatización del proceso.
    \item Cumplimiento normativo.
\end{itemize}

% ------------------------------------------------
\chapter{Marco Teórico}

\section{Flujo Operativo Actual de la Defensoría}

Descripción del proceso actual de recepción, clasificación y resolución de quejas dentro del IPN.

\section{Clasificación de Quejas}

Tipologías conforme al Reglamento Interno del IPN \cite{ref9}.

\section{Normativa Aplicable}

\begin{itemize}
    \item Reglamento Interno del IPN \cite{ref9}.
    \item Reglamentos internos de unidades académicas \cite{ref10}.
    \item Ley Federal de Protección de Datos Personales \cite{ref11}.
\end{itemize}

\section{Modelo de Clasificación Inteligente}

El sistema empleará aprendizaje supervisado y procesamiento de lenguaje natural para categorizar denuncias, incorporando revisión humana para mitigación de sesgos.

% ------------------------------------------------
\chapter{Metodología}

Se empleará la metodología Scrum-ban, combinando iteraciones planificadas (Scrum) con gestión visual del flujo de trabajo (Kanban), permitiendo adaptabilidad y mejora continua.

% ------------------------------------------------
\chapter{Conclusiones}

La plataforma propuesta representa una solución tecnológica orientada a mejorar la eficiencia, transparencia y seguridad en la gestión de denuncias dentro del Instituto Politécnico Nacional.

% ------------------------------------------------
\begin{thebibliography}{11}

\bibitem{ref1}
INEGI, Encuesta Nacional sobre Disponibilidad y Uso de Tecnologías de la Información en los Hogares (ENDUTIH) 2024.

\bibitem{ref2}
INEGI, Módulo sobre Ciberacoso (MOCIBA) 2024.

\bibitem{ref9}
Instituto Politécnico Nacional, Reglamento Interno del IPN, 2018.

\bibitem{ref10}
Instituto Politécnico Nacional, Reglamento General de Unidades Académicas, 2018.

\bibitem{ref11}
Ley Federal de Protección de Datos Personales en Posesión de los Particulares, 2018.

\end{thebibliography}

\end{document}
